% WHAT A WORD IS BROKEN WITH.
%
% \hyphenchar names the character a hyphenated word is broken with. Three
% cases, measured against pdftex with trip.tfm:
%
%   \hyphenchar\rip=`?    the font has no `?': the word is still broken,
%                         with NOTHING in front of the break, and
%                         \tracinglostchars names the character that is
%                         missing, as it names any other;
%   \hyphenchar\rip=`-    the font has `-': the usual break;
%   \hyphenchar\rip=-1    no hyphen character at all: the word is not
%                         broken.
%
% This engine left a word whose hyphen character the font had not got
% whole, and said nothing about it. trip breaks BBBBBB in a font whose
% \hyphenchar is `?'.
%
% Found and left open: the reference names the missing character while it
% is breaking the line that wants it, so under \tracingparagraphs the line
% falls among the trace; HSTeX breaks the words of a paragraph before it
% weighs any of them, so its line comes first.
\catcode`\{=1 \catcode`\}=2
\tracingonline=1 \tracinglostchars=1
\showboxdepth=99 \showboxbreadth=99
\output={\shipout\box255}
\font\rip=trip \rip
\lccode`A=1 \lccode`B=2 \lccode`C=3 \uchyph=1
\patterns{A1B B1B B1C C1A A1A}
\vsize=20000pt \hsize=10pt \parindent=0pt
\pretolerance=-1 \tolerance=10000 \hbadness=10000
\message{[A a hyphen character the font has not got]}
\hyphenchar\rip=`?
\setbox1=\vbox{\noindent A BBBB\par}
\showbox1
\message{[B one it has]}
\hyphenchar\rip=`-
\setbox1=\vbox{\noindent A BBBB\par}
\showbox1
\message{[C and none at all]}
\hyphenchar\rip=-1
\setbox1=\vbox{\noindent A BBBB\par}
\showbox1
\message{[done]}
\end
